\chapter{算法设计}\label{chap:algorithm}

\section{BFS 算法}

广度优先搜索（BFS）按照层次逐层扩展节点。在单位代价栅格地图中，BFS 能保证第一次到达终点时得到最短路径。其优点是实现简单、结果稳定；缺点是搜索具有较强盲目性，在大场景中会访问大量无关节点。

\section{A* 算法}

A* 在一致代价搜索基础上加入启发式函数，评价函数定义为
\begin{equation}
  f(n)=g(n)+h(n),
\end{equation}
其中 $g(n)$ 表示起点到当前节点的实际代价，$h(n)$ 表示当前节点到终点的估计代价。本文选用曼哈顿距离作为启发式函数：
\begin{equation}
  h(n)=|r-r_g|+|c-c_g|.
\end{equation}

相较于 BFS，A* 的搜索方向更集中，因此通常访问节点更少。

\section{改进 A* 算法}

基础 A* 只关注路径长度，不考虑路径是否贴近障碍，也不考虑转弯是否过多。为此，本文在单步代价中加入贴障惩罚和转向惩罚：
\begin{equation}
  c(n_i,n_{i+1})=1+\lambda_1 p_{\text{obs}}+\lambda_2 p_{\text{turn}}.
\end{equation}

这种设计的目标不是一味追求最短长度，而是在路径长度、平滑性和安全性之间取得更平衡的结果。

\section{Q-Learning 算法}

Q-Learning 属于强化学习方法。它通过不断与环境交互，更新状态-动作价值函数：
\begin{equation}
  Q(s,a)\leftarrow Q(s,a)+\alpha \left[r+\gamma \max_{a'}Q(s',a')-Q(s,a)\right].
\end{equation}

本文中，状态为机器人当前位置，动作是四个移动方向。奖励函数同时考虑了到达终点奖励、撞墙惩罚、靠近终点奖励和重复访问惩罚，以提高训练效率。

\section{人工势场法}

人工势场法将终点看作吸引源，将障碍物看作排斥源，机器人在合成势场作用下向目标移动。它的优点是实现直观、局部反应速度快，但缺点也很明显，即在复杂障碍环境中容易陷入局部极小值。

\section{Dijkstra 算法}

Dijkstra 是求解非负权图最短路径的经典方法。在本文的单位代价栅格环境中，它可以看作不带启发式信息的一致代价搜索。它能够稳定得到最短路径，但在复杂场景中搜索效率通常低于 A*。

\section{算法比较}

不同算法的特点如\cref{tab:method-comparison} 所示。

\begin{table}[!htb]
  \centering
  \caption{各算法特点对比}\label{tab:method-comparison}
  \begin{tabular}{@{}p{2.3cm}p{2.6cm}p{3.3cm}p{4.6cm}@{}}
    \toprule
    \textbf{算法} & \textbf{是否保证最短路} & \textbf{主要优点} & \textbf{主要缺点} \\
    \midrule
    BFS & 是 & 实现简单，结果稳定 & 搜索盲目性强，复杂场景扩展节点多 \\
    Dijkstra & 是 & 适用于非负权图，理论清晰 & 不利用目标方向信息，效率偏低 \\
    A* & 是 & 搜索效率较高，适合静态已知地图 & 仅以距离为目标时，路径可能贴近障碍 \\
    改进 A* & 综合代价更优 & 路径更平滑、更安全 & 参数需要调节，未必扩展最少节点 \\
    Q-Learning & 训练充分时可获得可行策略 & 可学习策略，适合扩展到智能体方法 & 训练时间较长 \\
    人工势场法 & 否 & 实现简单，局部反应快 & 易陷入局部极小值 \\
    \bottomrule
  \end{tabular}
\end{table}
